\documentclass{beamer}
% Оформление из предоставленного шаблона; содержание полностью заменено.
\beamertemplatenavigationsymbolsempty
\usecolortheme{beaver}
\setbeamertemplate{blocks}[rounded=true,shadow=true]
\setbeamertemplate{footline}[page number]
\setbeamertemplate{itemize item}{\textendash}
\setbeamertemplate{itemize subitem}{\textendash}
\setbeamertemplate{itemize subsubitem}{\textendash}
\usepackage{iftex}
\ifPDFTeX
  \usepackage[T2A]{fontenc}
  \usepackage[utf8]{inputenc}
\else
  \usepackage{fontspec}
  \setmainfont{cmunrm.otf}[BoldFont=cmunbx.otf,ItalicFont=cmunti.otf,BoldItalicFont=cmunbi.otf]
  \setsansfont{cmunss.otf}[BoldFont=cmunsx.otf,ItalicFont=cmunsi.otf,BoldItalicFont=cmunso.otf]
\fi
\usepackage[english,russian]{babel}
\usepackage{amssymb,amsfonts,amsmath}
\usepackage{booktabs,array}
\usepackage{tikz}
\usetikzlibrary{arrows.meta,positioning,calc}

\newcommand{\U}{\mathbf{u}}
\newcommand{\Grad}{\operatorname{grad}}
\newcommand{\Div}{\operatorname{div}}
\newcommand{\src}[1]{\par\vspace{0.5em}{\scriptsize\color{gray}#1\par}}
\definecolor{fieldblue}{HTML}{D8EAF9}
\definecolor{paramorange}{HTML}{FFE0B3}
\definecolor{opgreen}{HTML}{D9EDD9}
\tikzset{dag/.style={draw=black!60,rounded corners=2pt,fill=opgreen,
  minimum height=0.65cm,minimum width=1.35cm,inner sep=3pt,
  align=center,font=\small},
  field/.style={dag,fill=fieldblue},parameter/.style={dag,fill=paramorange},
  edge/.style={-{Stealth[length=2mm]},draw=black!75,line width=0.55pt}}

\title{Уравнения Навье--Стокса\protect\\в аэродинамике}
\subtitle{Лабораторная работа №\,1: вычислительный эксперимент}
\author{}
\institute{Foundation models for spatial-time series}
\date{10 октября 2026}


% Общий список рёбер: вход/операция. Используется и в дереве, и в матрице.
\newcommand{\TreeEdges}{1/2,2/4,3/4,4/5,6/7,7/9,8/9,9/10,5/11,10/11,13/14,12/15,14/15,11/16,15/16}
\newcommand{\PaperSource}{\href{https://doi.org/10.1103/9pwb-ky8n}{Hillestad et al., Physical Review Fluids 11, 014701 (2026)}}
\newcommand{\DataSource}{\href{https://doi.org/10.18710/SX0ZQT}{DataverseNO: doi:10.18710/SX0ZQT, V1 (2026)}}

\begin{document}

% 1
\begin{frame}
\thispagestyle{empty}
\maketitle
\vspace{-0.6cm}
\begin{center}\small
Вихревой след за моделью ветротурбины:\\
эксперимент Hillestad et al. (2026)\\
численное решение и подбор вязкости
\end{center}
\end{frame}

% 2
\begin{frame}{Базовое уравнение Навье--Стокса}
\begin{block}{Несжимаемая ньютоновская жидкость}
\[
\frac{\partial\U}{\partial t}+(\U\cdot\nabla)\U
=-\frac{1}{\rho}\nabla p+\nu\Delta\U,
\qquad \nabla\cdot\U=0.
\]
\end{block}
\[
\U=\U(x,y,z,t)=(u,v,w),\qquad
\mathbf{x}\in\Omega\subset\mathbb{R}^{3}.
\]
\small
Плотность $\rho$ и кинематическая вязкость $\nu$ постоянны.
Потенциал силы тяжести включён в модифицированное давление $p$.
Система координат неподвижна относительно лаборатории.

\medskip
В выбранном эксперименте рабочая среда --- вода.
Для решения нужны начальное поле скорости и граничные условия.
\end{frame}

% 3
\begin{frame}{Библиотеки и хранение дерева}
\begin{block}{NetworkX + SymPy}
\small
\texttt{NetworkX}: направленное дерево, атрибуты вершин,
проверка ацикличности и матрица смежности.

\medskip
\texttt{SymPy}: символические выражения, производные
и проверка тождеств.
\end{block}
\small
Для вершины сохраняются:
\begin{itemize}
\item идентификатор и операция;
\item тип результата: скаляр, вектор или матрица;
\item упорядоченный список аргументов \texttt{args}.
\end{itemize}
Каждое вхождение поля получает отдельную вершину:
получается дерево, являющееся частным случаем DAG.
На следующем слайде построена правая часть уравнения скорости.
\src{\href{https://networkx.org/documentation/stable/}{Документация NetworkX};
\href{https://docs.sympy.org/latest/}{документация SymPy}.}
\end{frame}

% 4
\begin{frame}{Дерево исходного уравнения Навье--Стокса}
\small
\[
\partial_t\U=F=-(\Grad\U)\U-\frac{\Grad p}{\rho}+\nu\Delta\U.
\]
\begin{center}
\begin{tikzpicture}[x=0.95cm,y=0.86cm,
 dag/.append style={minimum width=1.12cm,minimum height=0.49cm,font=\small,inner sep=2pt}]
\node[dag] (n16) at (0,0) {$16:\;+$};
\node[dag] (n11) at (-2.2,-1) {$11:\;+$};
\node[dag] (n15) at (3.6,-1) {$15:\;\times_s$};
\node[dag] (n5) at (-4.1,-2) {$5:\;-$};
\node[dag] (n10) at (-0.35,-2) {$10:\;-$};
\node[parameter] (n12) at (2.5,-2) {$12:\;\nu$};
\node[dag] (n14) at (4.7,-2) {$14:\;\Delta$};
\node[dag] (n4) at (-4.1,-3) {$4:\;\mathrm{mv}$};
\node[dag] (n9) at (-0.35,-3) {$9:\;\div_s$};
\node[field] (n13) at (4.7,-3) {$13:\;\U$};
\node[dag] (n2) at (-5,-4) {$2:\;\Grad$};
\node[field] (n3) at (-3.2,-4) {$3:\;\U$};
\node[dag] (n7) at (-1.3,-4) {$7:\;\Grad$};
\node[parameter] (n8) at (0.6,-4) {$8:\;\rho$};
\node[field] (n1) at (-5,-5) {$1:\;\U$};
\node[field] (n6) at (-1.3,-5) {$6:\;p$};
\foreach \s/\t in \TreeEdges {\draw[edge] (n\s)--(n\t);}
\end{tikzpicture}
\end{center}
\vspace{-0.3em}
\footnotesize
$(\Grad\U)_{ij}=\partial_j u_i$, поэтому $(\Grad\U)\U=(\U\cdot\nabla)\U$.
$\mathrm{mv}$ --- матрица на вектор; $\div_s$ --- деление вектора на скаляр;
$\times_s$ --- умножение на скаляр; $-$ --- смена знака.
\src{16 вершин, 15 рёбер. Стрелки: вход $\to$ операция.
Условие $\Div\U=0$ учитывается солвером отдельно.}
\end{frame}

% 5
\begin{frame}{Матрица смежности этого дерева}
\small
$Z_{ij}=1$, если вершина $j$ --- вход операции в вершине $i$.
\begin{columns}[c,onlytextwidth]
\begin{column}{0.53\textwidth}
\centering
\begin{tikzpicture}[x=0.27cm,y=0.27cm]
\node[font=\scriptsize] at (8.5,1.5) {$j$: вход};
\foreach \j in {1,...,16} {
 \node[font=\tiny] at (\j,0.6) {\j};
 \node[font=\tiny] at (0,-\j+0.5) {\j};
 \foreach \i in {1,...,16} {
  \filldraw[fill=opgreen!35!white,draw=white,line width=0.35pt]
    (\j-0.5,-\i+1) rectangle ++(1,-1);
  \node[font=\tiny,text=black!40] at (\j,-\i+0.5) {0};
 }
}
\foreach \s/\t in \TreeEdges {
 \filldraw[fill=black!70,draw=white,line width=0.35pt]
   (\s-0.5,-\t+1) rectangle ++(1,-1);
 \node[font=\tiny,text=white] at (\s,-\t+0.5) {1};
}
\node[font=\scriptsize] at (8.5,-16.5) {$i$: операция (строка)};
\end{tikzpicture}
\end{column}
\begin{column}{0.45\textwidth}
\footnotesize
{\scriptsize
\begin{tabular}{@{}rl@{}}
\toprule
Вершина & Содержимое \\
\midrule
1, 3, 13 & Поле $\U$ \\
2 & Матрица $\Grad\U$ \\
4, 5 & $(\Grad\U)\U$, затем минус \\
6, 7 & Поле $p$, затем $\Grad p$ \\
8 & Плотность $\rho$ \\
9, 10 & $\Grad p/\rho$, затем минус \\
11 & Сумма результатов 5 и 10 \\
12 & Вязкость $\nu$ \\
14, 15 & $\Delta\U$, затем $\nu\Delta\U$ \\
16 & Итоговая сумма $F$ \\
\bottomrule
\end{tabular}
}
\medskip

Топологический порядок вершин:
все входы предшествуют операции.
\end{column}
\end{columns}
\medskip
\footnotesize
В NetworkX при рёбрах «вход $\to$ операция» матрица
$Z=A^{\mathsf T}$. Порядок аргументов хранится отдельно:
$\texttt{args}(4)=(2,3)$, $\texttt{args}(9)=(7,8)$.
\end{frame}

% 6
\begin{frame}{Физическая постановка эксперимента}
\small
\begin{enumerate}
\item В водном канале NTNU активная решётка создаёт турбулентный
набегающий поток.
\item Поток приводит во вращение свободно вращающийся ротор
с двумя лопастями. Ротор --- лопасти вместе со ступицей.
\item За ротором, вниз по потоку, отслеживаются частицы-трассеры
и восстанавливается скорость воды.
\end{enumerate}
\medskip
Координата $x$ направлена вдоль основного потока;
$y,z$ задают поперечные направления. Начало $x$ --- плоскость ротора.
\begin{block}{Связь с аэродинамикой}
Это гидродинамическая модель ветротурбины.
Перенос выводов на течение воздуха требует сопоставления
$Re$, быстроходности и геометрии.
\end{block}
\src{Статья: \PaperSource. Данные установки: README архива.}
\end{frame}

% 7
\begin{frame}[shrink=5]{Экспериментальные данные для расчёта}
\small
\begin{block}{Какие данные доступны}
\begin{itemize}
\item Трёхмерные поля скорости $\U=(u,v,w)$ в последовательные моменты времени.
\item Пространственные координаты точек измерения и временная шкала.
\item Угол поворота ротора, характерные длина и скорость потока.
\end{itemize}
\end{block}
\begin{block}{Какие данные нужны для нашей задачи}
\begin{itemize}
\item Первое поле скорости и скорости на границе во всём временном окне
--- для начальных и граничных условий.
\item Скорости во внутренних точках --- для подбора вязкости
и проверки численного решения.
\end{itemize}
\end{block}
\footnotesize
Давление не измерено; оно определяется при решении уравнения.
Скорости получены методом PTV/Shake-the-Box и перенесены на сетку методом VIC\#.
\src{\DataSource.}
\end{frame}

% 8
\begin{frame}{Задача: решение Навье--Стокса и подбор вязкости}
\small
\textbf{Цель:} численно воспроизвести измеренное течение в трёхмерной
подобласти следа и подобрать кинематическую вязкость $\nu$.
\begin{block}{Прямая задача при заданном $\nu>0$}
\[
\begin{cases}
\partial_t\U_\nu+(\U_\nu\cdot\nabla)\U_\nu
=-\rho^{-1}\nabla p_\nu+\nu\Delta\U_\nu,\\[0.3em]
\nabla\cdot\U_\nu=0,
\end{cases}
\quad (\mathbf x,t)\in\Omega\times[t_0,t_1].
\]
\end{block}
\begin{itemize}
\item $\Omega$ --- часть измеряемого объёма, не пересекающая лопасти и опору.
Выбрано: $x\in[20,100]$ мм, $y,z\in[-24,24]$ мм.
\item Неизвестные поля: скорость $\U_\nu$ и давление $p_\nu$.
Плотность $\rho$ постоянна.
\item При каждом пробном $\nu$ решается одна и та же начально-краевая задача;
её решение сравнивается с экспериментом.
\end{itemize}
\footnotesize
Это постановка лабораторной по данным случая 1L;
авторы статьи исследовали динамику вихрей, а не подбирали $\nu$ этим способом.
\end{frame}

% 9
\begin{frame}{Начальные и граничные условия}
\small
\begin{block}{Данные задают условия задачи}
\[
\U_\nu(\mathbf x,t_0)=\U_0(\mathbf x),\qquad
\U_\nu|_{\partial\Omega}=\mathbf g(\mathbf x,t).
\]
\end{block}
\begin{itemize}
\item $\U_0$ получается из первого кадра окна;
$\mathbf g$ --- из измерений на границе для всего окна.
Они одинаковы при всех пробных $\nu$.
\item Перед расчётом: перевод мм в м, интерполяция на расчётную сетку
и согласование данных с несжимаемостью.
\end{itemize}
\[
\nabla\cdot\U_0=0,\qquad
\U_0|_{\partial\Omega}=\mathbf g(\cdot,t_0),\qquad
\int_{\partial\Omega}\mathbf g\cdot\mathbf n\,dS=0.
\]
Давление вычисляется солвером. При заданной скорости на всей границе
его постоянный уровень фиксируется условием
\[
\int_\Omega p_\nu(\mathbf x,t)\,dV=0.
\]
\footnotesize
Границы подобласти проходят в жидкости; условие нулевой скорости на них
не применяется. Коррекция измерений должна согласовать начальные и граничные поля.
\end{frame}

% 10
\begin{frame}{Обратная задача и оптимизатор}
\small
Для сравнения используются точки $\mathbf x_j\in\Omega_{\mathrm{int}}$,
лежащие внутри объёма. $H_j$ переносит расчётное поле в точку измерения.
\[
\begin{aligned}
\nu^*&=\arg\min_{\nu>0}J(\nu),\\[0.3em]
J(\nu)&=\sum_{n\in\mathcal T_{\mathrm{fit}}}\sum_j
\left\|\frac{H_j\U_\nu(t_n)-\U_{\mathrm{exp}}(\mathbf x_j,t_n)}{U_\infty}\right\|^2.
\end{aligned}
\]
\begin{block}{Выбор: SciPy least\_squares, метод TRF}
Оптимизатор получает вектор разностей скоростей.
Каждая оценка целевой функции требует численного решения PDE.
\end{block}
\[
\nu=\nu_{\mathrm{ref}}e^\eta,\qquad
\eta\ \longrightarrow\ \nu\ \longrightarrow\
\operatorname{SolveNS}(\nu;\U_0,\mathbf g)\ \longrightarrow\ J(\nu).
\]
\footnotesize
Оптимизируется один безразмерный параметр $\eta$; $\nu_{\mathrm{ref}}>0$ ---
опорная вязкость. Поиск: $\nu/\nu_{\mathrm{ref}}\in[0.5,2]$. Это исследовательский диапазон, без измеренной температуры воды.
Граничные точки и начальный кадр не входят в функцию ошибки.
\src{\href{https://docs.scipy.org/doc/scipy/reference/generated/scipy.optimize.least_squares.html}{SciPy: least\_squares и Trust Region Reflective (TRF)}.}
\end{frame}

% 11
\begin{frame}[shrink=5]{Реализованный солвер: трёхмерная сетка MAC}
\small
\textbf{Пространство:} скорости на гранях ячеек, давление в центрах.
Основная сетка: $36\times24\times24$; шаг около $2$ мм.
Перенос --- противопоточная схема первого порядка;
диффузия --- центральная схема второго порядка.
\medskip

\textbf{Время:} SSP-RK3 с проекцией на каждой стадии:
\[
\Delta_h q=\frac{\nabla_h\cdot\U^*}{\alpha\Delta t},\qquad
\U^{\mathrm{new}}=\U^*-\alpha\Delta t\nabla_h q.
\]
\begin{itemize}
\item Задача Пуассона: DCT-II, нулевой нормальный градиент поправки,
нулевое среднее; границы согласованы по потоку.
\item Адаптивный шаг: ограничение переноса CFL $0.35$,
диффузии и ближайшего момента измерения.
\item Реализация: NumPy/SciPy, файл \texttt{lab01/solver.py}.
\end{itemize}
\footnotesize
Центральный перенос оказался неустойчивым на сгущённой сетке.
Противопоточная схема устойчива, но вносит численную диффузию.
FEniCSx отсутствует в окружении; фактический расчёт выполнен MAC-солвером.
\src{\href{https://doi.org/10.1063/1.1761178}{Harlow \& Welch (1965): MAC};
\href{https://doi.org/10.1090/S0025-5718-98-00913-2}{Gottlieb \& Shu (1998): SSP-RK}.}
\end{frame}

% 12
\begin{frame}[shrink=5]{Аппроксимация, ошибка и графики сравнения}
\small
Аппроксимация данных --- поле $\U_{\nu^*}$, полученное численным
решением PDE при подобранной вязкости.
\[
E_{\mathrm{val}}=
\left[
\frac{\displaystyle\sum_{n\in\mathcal T_{\mathrm{val}}}\sum_j
 \|H_j\U_{\nu^*}(t_n)-\U_{\mathrm{exp}}(\mathbf x_j,t_n)\|^2}
{\displaystyle\sum_{n\in\mathcal T_{\mathrm{val}}}\sum_j
 \|\U_{\mathrm{exp}}(\mathbf x_j,t_n)\|^2}
\right]^{1/2}.
\]
Подбор: $t=0$--$0.09$ с; контроль: $t=0.10$--$0.17$ с.
Окна не пересекаются. В контрольном окне $\nu^*$
фиксировано; начальное поле и границы задаются по его измерениям.
\medskip

\textbf{Графики построены в ноутбуке}
\begin{enumerate}
\item Компонента скорости во времени: измерения и расчёт на одной оси.
\item Одно сечение объёма: поле измерений, расчёт и карта ошибки.
\item $J(\nu)$: выраженность минимума и чувствительность к вязкости.
\end{enumerate}
\footnotesize
Начальный кадр исключается из оценки. При известных измеренных границах
проверяется воспроизведение течения; самостоятельный прогноз здесь не заявляется.
\end{frame}

% 13
\begin{frame}[shrink=5]{Параметры и безразмерное уравнение}
\small
Масштабы: $x^*=x/D$, $t^*=tU_\infty/D$,
$\U^*=\U/U_\infty$, $p^*=p/(\rho U_\infty^2)$.
\begin{block}{Исходное уравнение в безразмерных переменных}
\[
\begin{gathered}
\partial_{t^*}\U^*+(\U^*\cdot\nabla^*)\U^*
=-\nabla^*p^*+\frac{1}{Re_D}\Delta^*\U^*,\\[0.3em]
\nabla^*\cdot\U^*=0.
\end{gathered}
\]
\end{block}
\[
\nu_{\mathrm{ref}}=\frac{U_\infty D}{Re_{D,\mathrm{ref}}}
\approx9\cdot10^{-7}\ \text{м}^2/\text{с},\qquad
\frac{D}{U_\infty}=0.1125\ \text{с}.
\]
\begin{itemize}
\item Подбирается $\nu$; коэффициент $1/Re_D=\nu/(U_\infty D)$ меняется вместе с ним.
\item $I$ и угол ротора характеризуют условия течения; они не заменяют $\nu$.
\item Малый $1/Re_D$ сам по себе не обосновывает удаление вязкого члена.
\end{itemize}
\src{Масштабы из README случая 1L. $Re_{D,\mathrm{ref}}\approx2\cdot10^4$ задаёт опорную оценку, а не результат подбора.}
\end{frame}

% 14
\begin{frame}{Надёжность определения вязкости}
\small
\begin{itemize}
\item При $Re_D\approx2\cdot10^4$ чувствительность скорости к $\nu$
оказалась слабой: размах $J(\nu)$ на всём диапазоне --- $0.135\%$.
Минимум на верхней границе; надёжная оценка вязкости не получена.
\item Выполнены расчёты на трёх сетках и с вдвое меньшим ограничением
шага по времени; сравниваются одинаковые 539 внутренних точек.
\item Шум, обработка измерений и численная диффузия могут сдвинуть
минимум. Подобранное число нельзя автоматически считать точной
молекулярной вязкостью воды.
\end{itemize}
\medskip
Проверки качества самого решения:
\[
\|\nabla\cdot\U_{\nu^*}\|_{L^2(\Omega)},\qquad
\left|\int_{\partial\Omega}\U_{\nu^*}\cdot\mathbf n\,dS\right|.
\]
\footnotesize
Параметр оценивается по мгновенным полям. При фильтрации или осреднении
нужно учитывать, что точные уравнения для такого поля содержат дополнительные напряжения.
\end{frame}

% 15
\begin{frame}{Вычислительный эксперимент выполнен}
\small
Загружены и проверены 18 кадров случая 1L: 2732, 2737, \ldots, 2817.
Шаг выбранных данных --- $0.01$ с, длительность --- $0.17$ с.
Полный архив не скачивался; сохранены CRC32 и SHA-256 файлов.
\medskip

\begin{tabular}{@{}lrr@{}}
\toprule
Сетка / CFL & Ошибка подбора & Ошибка контроля \\
\midrule
$24\times16\times16$ / $0.35$ & $25.83\%$ & $24.75\%$ \\
$36\times24\times24$ / $0.35$ & $25.50\%$ & $24.46\%$ \\
$48\times32\times32$ / $0.35$ & $25.35\%$ & $24.36\%$ \\
$36\times24\times24$ / $0.175$ & $25.50\%$ & $24.46\%$ \\
\bottomrule
\end{tabular}
\medskip

\begin{itemize}
\item Подбор TRF: $\nu^*=1.8\cdot10^{-6}$ м$^2$/с, верхняя граница поиска.
\item Основная сетка: контрольный RMSE $0.03318$ м/с;
максимум RMS дискретной дивергенции $8.8\cdot10^{-15}$ с$^{-1}$.
\item Построены временные кривые, сечение поля, карта ошибки и $J(\nu)$.
\end{itemize}
\footnotesize
Контроль использует собственное начальное поле и известные измеренные границы.
Численный эксперимент выполнен; точная молекулярная вязкость не определена.
\end{frame}

% 16
\begin{frame}{Источники}
\small
\begin{enumerate}
\item Hillestad J. N. et al.
\emph{Volumetric three-dimensional experimental measurement of
vortex dynamics in a rotating wake}.
Physical Review Fluids \textbf{11}, 014701 (2026).\\
\href{https://doi.org/10.1103/9pwb-ky8n}{doi:10.1103/9pwb-ky8n}.
\smallskip
\item Hillestad et al. \emph{Volumetric three-dimensional experimental
measurement of a model wind turbine wake}.
DataverseNO, V1 (2026), случай 1L и \texttt{00\_README.txt}.\\
\href{https://doi.org/10.18710/SX0ZQT}{doi:10.18710/SX0ZQT}.
\smallskip
\item Harlow F. H., Welch J. E. (1965): разнесённая сетка MAC.\\
\href{https://doi.org/10.1063/1.1761178}{doi:10.1063/1.1761178}.\\
Gottlieb S., Shu C.-W. (1998): SSP/TVD Runge--Kutta.\\
\href{https://doi.org/10.1090/S0025-5718-98-00913-2}{doi:10.1090/S0025-5718-98-00913-2}.
\smallskip
\item SciPy. \emph{scipy.optimize.least\_squares}:
нелинейный метод наименьших квадратов, алгоритм TRF.\\
\href{https://docs.scipy.org/doc/scipy/reference/generated/scipy.optimize.least_squares.html}{Документация scipy.optimize.least\_squares}.
\end{enumerate}
\end{frame}

\end{document}
